\documentclass{IOS-Book-Article}

% -----------------------------------------------------------------------------
% Packages
% -----------------------------------------------------------------------------
\usepackage{mathptmx}
\usepackage{soul}
\setuldepth{article}
\usepackage{graphicx}
\usepackage{url}
\usepackage{booktabs}
\usepackage{array}
\usepackage{float}

% IOS Press template helper.
\def\hb{\hbox to 11.5 cm{}}

% -----------------------------------------------------------------------------
% Document
% -----------------------------------------------------------------------------
\begin{document}

\begin{frontmatter}

\title{Lextract: LLM-Based Extraction of Relevant Market Definitions from
European Commission Decisions}
\runningtitle{Lextract: Market Definition Extraction}

\author{%
  \fnms{Shriyan S.} \snm{Yamali}%
  \orcid{0009-0002-4861-7264}%
}
\address{Independent Researcher, United States}
\runningauthor{S.S. Yamali}

\begin{abstract}
Relevant market definitions are central to merger and antitrust analysis, but
collecting them at scale from European Commission decisions is difficult because
the decisions are numerous and heterogeneous in structure and wording. We present
Lextract, an open-source pipeline that filters decision PDFs, uses Gemini 2.0
Flash in two calls to locate the market-definition section and to isolate
individual definitions, and stores the results as JSON. Applied to 600 cases, it
produced 4,105 records that feed JurisMercatus, a semantic-search interface. We
document the prompts, settings, exclusion rules and thresholds, give a worked
example from source text to record, and combine the model with three symbolic
quality checks, namely a structural validity test derived from 50 manually
checked definitions, a source-fidelity test that checks whether each record
occurs verbatim in its decision, and a manual inspection of 100 records. Of the
records, 76.6\% pass the structural test. Against all 699 retrievable
decisions, 79\% of the 4,067 verifiable records are verbatim and 9\%
near-verbatim. The fidelity test flags the remaining 12\% as absent from their
source, three quarters of them from short simplified-procedure decisions without
market analysis, and removes them in a cleaned release of 3,607 records. The
structural test accepts 71\% of the flagged records, which shows that checking
output against the source text adds information that resemblance-based checks
miss. Annotation sheets and baseline scripts for a gold-standard evaluation are
provided.
\end{abstract}

\begin{keyword}
legal information extraction\sep
competition law\sep
relevant market definition\sep
large language models\sep
quality control\sep
European Commission
\end{keyword}

\end{frontmatter}

% -----------------------------------------------------------------------------
\section{Introduction}
\label{sec:introduction}
% -----------------------------------------------------------------------------

Relevant market definition is a basic step in merger and antitrust analysis
because it fixes the products and the geographic area within which market power
is assessed \cite{marketnotice2024}. Practitioners and researchers therefore
need to know how the European Commission has delineated a given market before,
for instance whether retail distribution of cars is national or EEA-wide,
whether a product market was segmented or ``left open'', and how such findings
changed over time. This is a concrete user problem. The Commission has
published more than 6,000 merger and antitrust decisions \cite{bldr2024} and
adds roughly 280 a year \cite{apdtsf2021}. The market-definition discussion sits
under different headings and in different wording, so keyword search over PDFs
is unreliable and manual collection is slow and hard to reproduce. The only
dedicated resource we know of, the Caselex Market Definitions Module
\cite{caselex}, is a paid, editorially curated web product without bulk or
programmatic access.

This paper presents \emph{Lextract},\footnote{\url{https://github.com/shriyanyamali/Lextract/tree/lda2026}}
an open-source pipeline that extracts market-definition passages from
Commission decision PDFs with a large language model (LLM) and stores them as
structured records, and the 4,105-record collection behind
\emph{JurisMercatus},\footnote{\url{https://jurismercatus.shriyanyamali.com}} a
semantic-search interface. The design follows the hybrid view that LLMs are
good at locating and copying legally relevant passages and that their output
should be checked by deterministic means. The paper makes four contributions.
\begin{itemize}
  \item A complete description of the pipeline, covering prompts, model and
        settings, exclusion phrases and thresholds, with a worked example from
        source text to record (Sections~\ref{sec:pipeline}--\ref{sec:example}).
  \item A description of the released collection and of what researchers can
        do with it (Section~\ref{sec:dataset}).
  \item A three-part evaluation combining a structural validity test, a
        \emph{source-fidelity} test that checks whether each record occurs
        verbatim in its decision, and a manual inspection of 100 records
        (Sections~\ref{sec:evaluation}--\ref{sec:results}).
  \item An analysis of the records the fidelity test flags, which traces them
        to short-form decisions and to example passages in the prompt, together
        with the corresponding fixes (Section~\ref{sec:errors}) and the tools
        for a gold-standard evaluation (Section~\ref{sec:limitations}).
\end{itemize}

% -----------------------------------------------------------------------------
\section{Related Work and Existing Resources}
\label{sec:related-work}
% -----------------------------------------------------------------------------

Legal information extraction combines deterministic, statistical and, more
recently, generative methods. Agnoloni et al.\ presented the BO-ECLI Parser for
extracting and linking legal references from European case law
\cite{agnoloni2017boecli}, and Darji et al.\ released annotated German
legal-reference data and stressed the role of domain-specific gold data
\cite{darji2023german}. Recent work uses LLMs for semantic structuring.
Janatian et al.\ turned legislation into structured pathways for expert systems
\cite{janatian2023text}; Costa et al.\ extracted elements from Brazilian
judicial decisions, evaluated them with precision and recall against expert
annotations and published them with provenance \cite{costa2024semantic}; Spada
et al.\ extracted events and dates from EU procurement notices with expert
validation \cite{spada2025expert}; others address data scarcity and context
\cite{nguyen2025data,bjbmmcmnthsktczmm2025,csfafemspsrdsc2024}. Lextract shares
their pipeline view, targets complete market-definition passages and
emphasises inexpensive deterministic checks of LLM output, which the
hallucination literature motivates \cite{ktanovssze2025}.

\paragraph{Relation to existing resources.} Caselex \cite{caselex} offers
human-reviewed market definitions across many authorities and languages through
a subscription and a web interface. Lextract complements it with a free,
reproducible and updatable collection that offers programmatic access. The
pipeline can be re-run on new decisions, every record carries case metadata and
the decision URL, and an export tool writes CSV, JSONL, Parquet, SQLite, BibTeX
and RIS. To our knowledge no openly licensed, machine-readable collection of
Commission market definitions existed before.

% -----------------------------------------------------------------------------
\section{The Lextract Pipeline}
\label{sec:pipeline}
% -----------------------------------------------------------------------------

Lextract is a sequence of six Python scripts driven by \texttt{run\_pipeline.py}.
The prompts are reproduced verbatim in the repository
(\texttt{docs/PROMPTS.md}) and summarised here.

\subsection{Document Collection and Filtering}

\emph{Step 1} reads the Excel export of the Commission's competition case-search
portal and extracts, per case, the decision URL, case number, year and policy
area, the last inferred from the URL path. \emph{Step 2} downloads each PDF and
extracts its text with PyPDF2. A document is excluded if it has exactly three
pages and contains one of five phrases from the closing recital of short-form
decisions adopted under the simplified procedure \cite{simplified2013}, namely
``For the reasons set out in the Notice on a simplified'', ``European Commission
has decided not to oppose the notified operation'', ``declare it compatible with
the internal market'', ``This decision is adopted in application of Article
6(1)(b)'' and ``Merger Regulation and Article 57 of the EEA Agreement''. Such
clearances contain no market analysis, so any ``definition'' extracted from them
would be spurious. Included documents are written to batch files with a
metadata header and split into two groups at 80,000 characters so that long
decisions can be processed separately; no other length filter is applied.
Version 1.1.0 of the code matches the phrases independently of whitespace and
applies the rule to documents of at most three pages (Section~\ref{sec:errors}).

\subsection{Section Identification (LLM call 1)}

\emph{Step 3} sends each batch file to Gemini 2.0 Flash
(\texttt{gemini-2.0-flash} via the \texttt{google-generativeai} SDK) with
default generation settings, that is, without temperature, top-$p$ or
output-length overrides, system instruction or JSON mode. The prompt asks the
model to find the first occurrence of the words ``market definition'', to start
one line before it, to continue past sub-headings and to stop at the next
heading of the same level (e.g.\ from ``IV. Market Definition'' to
``V.~\ldots''), excluding that heading, and to reproduce the metadata header
unchanged. The output is the market-definition section as plain text.

\subsection{Definition Extraction (LLM call 2)}

\emph{Step 4} sends each section to the same model with a prompt that asks for a
JSON array of objects with the fields \texttt{case\_number}, \texttt{year},
\texttt{policy\_area}, \texttt{link}, \texttt{topic} and \texttt{text}. It
instructs the model to extract ``the entirety of each individual market
definition'' as ``a direct quote of the multiple paragraphs'', to treat ``left
open'' passages as definitions, to drop paragraph numbers, footnote markers and
footnotes, to re-join words split by PDF extraction and to assign a short
\texttt{topic}. It states that a definition ``is always at least 4 sentences
long, but never longer than 12 sentences'' and contains four worked examples on
detergents, chemical intermediates, flame retardants and trimethylolpropane.
The two-stage design limits the second call to a few pages of relevant text
instead of a full decision.

\subsection{Cleaning and Merging}

\emph{Step 5} strips Markdown code fences from the model's output and
\emph{Step 6} concatenates all files into \texttt{data/output.json}. No
de-duplication, truncation or editing of the text is performed, so the released
records are exactly what the model returned.

% -----------------------------------------------------------------------------
\section{Worked Example}
\label{sec:example}
% -----------------------------------------------------------------------------

\begin{table}[t]
\caption{Worked example from Case M.9324. The upper part shows the PyPDF2 text
of the decision, with footnotes and later paragraphs omitted, and the lower part
shows the record produced by Lextract.}
\label{tab:example}
\small
\begin{tabular}{p{0.97\textwidth}}
\toprule
\textbf{Source text (decision PDF, Section 4)}\\
\texttt{4. PRODUCT MARKET DEFINITION}\\
\texttt{4.1.1. Wholesale distribution of ICT and CE products}\\
\texttt{(12) In the precedents listed at paragraphs (13)-(15), the Commission assessed the possible product markets for wholesale distribution of IT products, telecommunication equipment and CE.}\\
\texttt{(13) In Avnet/Magirus,8 Tech Data/Scribona,9 Tech Data/Avnet,10 the Commission has defined a market for the wholesale distribution of IT products, but left open the question whether the market should be further sub -segmented (i) by product categories (e.g. server, storage, PCs, etc.), (ii) by direct and indirect sales11 or (iii) by different distribution models (broadline vs value -added distribution ("VAD")).12}\\
\texttt{(14) In Argues/SHC13 and ALSO/PCF,14 the Commission identified a separate market [\ldots]}\\
\midrule
\textbf{Extracted record}\\
\texttt{\{"case\_number": "M.9324", "year": "2019", "policy\_area": "Merger",}\\
\texttt{\ "link": "}\url{https://ec.europa.eu/competition/mergers/cases1/202048/m9324_479_3.pdf}\texttt{",}\\
\texttt{\ "topic": "Wholesale distribution of IT products",}\\
\texttt{\ "text": "In Avnet/Magirus, Tech Data/Scribona , Tech Data/Avnet , the Commission has defined a market for the wholesale distribution of IT products, but left open the question whether the market should be further sub -segmented (i) by product categories (e.g. server, storage, PCs, etc.), (ii) by direct and indirect sales or (iii) by different distribution models (broadline vs value -added distribution ("VAD"))."\}}\\
\bottomrule
\end{tabular}
\end{table}

Table~\ref{tab:example} follows one passage through the pipeline. The source is
Case M.9324 \emph{ALSO / ABC Data Group} (12 June 2019), a 16-page Article
6(1)(b) decision. Call 1 located the heading ``4. PRODUCT MARKET DEFINITION''
and returned Section 4 up to, but excluding, ``5. COMPETITIVE ASSESSMENT''.
Call 2 returned seven records, and the one shown is paragraph (13). The model
removed the paragraph number and footnote markers 8--12, kept the wording
otherwise intact, including the extraction artefact ``sub -segmented'', and
labelled the topic. The structural test accepts the record (63 words, learned
phrase ``the commission has'', marker ``left open''), and the fidelity test
finds it verbatim in the decision (coverage 1.0). The sibling record for
paragraph (14) is near-verbatim (coverage 0.83) because in the PyPDF2 text that
paragraph is interrupted by a page-break footnote block, which the model
correctly skipped. The other five records for this decision are verbatim.

% -----------------------------------------------------------------------------
\section{The JurisMercatus Collection}
\label{sec:dataset}
% -----------------------------------------------------------------------------

The collection (\texttt{data/jurismercatus\_definitions.json} in the
repository) contains 4,105 records from 600 cases (578 merger, 22 antitrust)
and 704 distinct decision documents dated 2016--2025 (Table~\ref{tab:years}).
The documents were exported in English from the case-search portal, and the
released links identify each processed document. Cases contribute 1 to 45
records (median 6), records range from 9 to 1,182 words (median 67, mean 91),
and the \texttt{topic} field holds 3,250 distinct free-text labels. A cleaned
release (\texttt{data/jurismercatus\_definitions\_clean.json}, 3,607 records,
455 cases) omits the records that the fidelity test (Section~\ref{sec:results})
finds absent from their source and all records from documents of at most three
pages.

\begin{table}[t]
\caption{Records per year and structural validity rate.}
\label{tab:years}
\small
\centering
\setlength{\tabcolsep}{4.5pt}
\begin{tabular}{lrrrrrrrrrrr}
\toprule
Year & 2016 & 2017 & 2018 & 2019 & 2020 & 2021 & 2022 & 2023 & 2024 & 2025 & All\\
\midrule
Records   & 9  & 66 & 621 & 492 & 699 & 594 & 691 & 618 & 297 & 18 & 4,105\\
Valid     & 3  & 45 & 467 & 421 & 551 & 477 & 509 & 472 & 184 & 14 & 3,143\\
Rate (\%) & 33 & 68 & 75  & 86  & 79  & 80  & 74  & 76  & 62  & 78 & 76.6\\
\bottomrule
\end{tabular}
\end{table}

\paragraph{What researchers can do with it.} Each record pairs a passage with
its case number, year, policy area and decision URL, so it can be cited and
checked. Through JurisMercatus, a query such as ``retail distribution of new
passenger cars, geographic scope'' returns the passages in which the Commission
discussed that question across cases, ranked by embedding similarity. Exported
to a spreadsheet or SQLite database, the records support questions such as how
often the Commission leaves a definition open by sector and year, or how
geographic findings (national, EEA-wide, worldwide) are distributed.

% -----------------------------------------------------------------------------
\section{Evaluation Design}
\label{sec:evaluation}
% -----------------------------------------------------------------------------

\subsection{Reference Set and Structural Validity Test}

We evaluate the collection with three complementary instruments, each of which
measures a different property. The reference set
(\texttt{evaluation/reference.json}) holds 50 records that the author selected
from the pipeline output and checked against the source PDFs. They cover 39
merger cases, mostly from 2022--2024, and range from 33 to 557 words (median
135).

The structural test (\texttt{evaluate.py}) derives three checks from the
reference set. The length check requires 26 to 668 words, the reference range
widened by 20\% on each side. The phrase check requires at least one of the 121
word $n$-grams ($n=2$--5) that occur in at least 15\% of reference records, such
as ``can be left open'' or ``for the purposes of this decision''. The marker
check requires at least one of 18 fixed markers, namely ``relevant market'',
``product market'', ``geographic market'', ``market definition'', ``left
open'', ``EEA-wide'', ``EEA wide'', ``national market'', ``Commission
considers'', ``Commission found'', ``Commission has previously'', ``Notifying
Party'', ``internal market'', ``serious doubts'', ``plausible market'', ``exact
scope'', ``for the purposes of'' and ``cannot be excluded''. Seven of the 50
reference records fail the marker check themselves, so the markers form a
strict proxy. The test measures whether a record resembles a verified
market-definition passage in length and phrasing. Whether the text occurs in
the decision is measured by the fidelity test, while passage boundaries and
missed definitions are left to the gold-standard study
(Section~\ref{sec:limitations}). We therefore report the result as a
\emph{structural validity rate}.

\subsection{Source-Fidelity Test}

The fidelity test (\texttt{verify\_fidelity.py}) asks whether a record is
copied from its decision. Record and PyPDF2 decision text are normalised to
lower-case letters only, which removes whitespace, digits and punctuation and so
neutralises paragraph numbers, footnote markers and split-word artefacts. The
record is cut into 40-letter windows at stride 20, and \emph{coverage} is the
fraction of windows found in the source. Records with coverage of at least 0.95
count as \emph{verbatim}, from 0.80 to 0.95 as \emph{near-verbatim}, from 0.50
to 0.80 as \emph{partial} and below 0.50 as \emph{absent}. Whether a copied
passage is the right one is left to the gold-standard study. The test was run in
October 2026 on every record whose decision could be fetched from the stored
URL, covering 699 of the 704 documents and 4,067 of the 4,105 records; four URLs
of case M.11241 now redirect and one points to an HTML index page.

\subsection{Manual Inspection}

To interpret the structural test's verdicts, 50 passing and 50 failing records
were drawn at random (seeded) and read. Each was labelled from the record text
alone. Label A1 marks a market-definition statement, such as a precedent, a
party submission, a market-investigation finding or a Commission conclusion on
product or geographic scope. Label A2 marks reasoning within a
market-definition analysis without a definitional claim, B a fragment such as a
bare heading, and C text outside market definition, such as competitive
assessment, remedies, party descriptions or legal background. The labels are
released with the sample (\texttt{evaluation/manual\_inspection\_sample.csv}).
This single-annotator reading checks face validity. Borderline cases include a
conclusion that a transaction raises no serious doubts ``regardless of the exact
product market definition'', labelled C, and a list of ``affected markets'',
labelled A1. Precision, recall and F1 require the full list of definition
passages in whole decisions, and the choice of annotation unit (paragraph,
sub-section or conclusion) affects the scores. The repository includes an
annotation protocol and a script that produces, for $N$ random decisions, a
sheet listing every record next to the decision text with columns for true
positive, boundary error, false positive and missed definition
(Section~\ref{sec:limitations}).

% -----------------------------------------------------------------------------
\section{Results}
\label{sec:results}
% -----------------------------------------------------------------------------

\subsection{Structural Validity}

Of the 4,105 records, 3,143 (76.6\%) pass all three checks. The 962 failures
divide into 668 marker only, 185 length only, 80 length and marker, 19 phrase
and marker, and 10 all three. Of the 275 length failures, 261 are shorter than
26 words and 14 longer than 668. The yearly rate (Table~\ref{tab:years}) lies
between 68\% and 86\% in every year with more than 50 records except 2024
(62\%), and the 65 antitrust records reach 62\%. Among the 777 records without a
marker, 535 (69\%) still contain the word ``market'', 105 ``national'' and 65
``EEA'', so many express scope in wording the marker list does not cover, such
as ``the market is national in scope''; 242 (31\%) do not mention a market.

\subsection{Source Fidelity}

\begin{table}[t]
\caption{Source-fidelity test on all retrievable records and on the reference
set's source documents (five URLs, 38 records, were unretrievable).}
\label{tab:fidelity}
\small
\centering
\begin{tabular}{lrr}
\toprule
 & Collection & Reference set\\
\midrule
Documents checked (of which $\le$3 pages) & 699 (136) & 38 (1)\\
Records checked & 4,067 & 48\\
Verbatim (coverage $\ge$0.95) & 3,221 (79.2\%) & 44\\
\quad of which exact match (1.0) & 2,985 (73.4\%) & --\\
Near-verbatim (0.80--0.95) & 349 (8.6\%) & 3\\
Partial (0.50--0.80) & 26 (0.6\%) & 0\\
Absent ($<$0.50) & 471 (11.6\%) & 1\\
Absent, from $\le$3-page documents & 353 of 380 & 1 of 1\\
Absent, from longer documents & 118 of 3,687 & 0 of 47\\
\bottomrule
\end{tabular}
\end{table}

Of the 4,067 verifiable records (Table~\ref{tab:fidelity}), 3,221 (79.2\%) are
verbatim, 349 (8.6\%) near-verbatim, 26 (0.6\%) partial and 471 (11.6\%)
absent. Absent records concentrate in documents without market analysis. Of the
699 documents, 136 have at most three pages, and 353 of their 380 records are
absent; the remaining 27 are copied boilerplate such as party descriptions. Of
the 118 absent records from longer documents, 26 reproduce the prompt's example
passages (Section~\ref{sec:errors}), 32 come from twelve documents of four to
seven pages, and 60 come from 19 longer documents. The latter are mainly
ancillary decisions that refer back to a main clearance decision, whose records
restate the main decision's definitions in words that do not occur in the
document, and a few long decisions whose records share only fragments with the
text. Of the 349 near-verbatim records, 292 have at most two missing windows,
which corresponds to a single break. Inspection of such cases in twelve
documents found two causes, footnote text interleaved at a page break in the
PyPDF2 text, as in the worked example, and a skipped intermediate paragraph,
such as a party's submission between the Commission's precedent and its
conclusion in M.10745. No paraphrase was found.

Of the 48 reference records whose documents could be retrieved, 44 are verbatim,
three near-verbatim and one absent. Of the 471 absent records, 336 (71\%) pass
the structural test, close to the rate for verbatim records (2,475 of 3,221,
77\%), so structural resemblance does not indicate presence in the source.

\subsection{Manual Inspection}

Among the 50 passing records, 49 were labelled A1 and one C, a
competitive-assessment conclusion. One of the 49 is a prompt example that the
fidelity test flags (Section~\ref{sec:errors}). Among the 50 failing records, 38
were labelled A1 and three A2, so 41 (82\%, interval 69--90\%) contain
market-definition content. One is a bare heading (B), and eight (16\%) fall
outside market definition (C), covering remedies reasoning, a party
description, Treaty text and competitive assessment. Seven of the 41 are
complete one-sentence scope statements below the 26-word minimum, such as ``For
AIs, the Commission therefore takes the view that the markets are at least
EEA-wide'', and four of them fail on length alone. One is a coherent 727-word
party submission that exceeds only the maximum. Most structural failures
therefore stem from wording and length, and fidelity is measured separately
(Section~\ref{sec:results}).

% -----------------------------------------------------------------------------
\section{Error Analysis and Fixes}
\label{sec:errors}
% -----------------------------------------------------------------------------

\paragraph{Short-form decisions.} Most absent records, 353 of 471, come from
simplified-procedure clearances. These contain the exclusion phrases, but PyPDF2
renders them with inserted spaces and line breaks, as in ``on a simpli fied
procedure'', so exact matching does not find them. In a 40-document pilot
sample, none of eight three-page documents matched a phrase, and a two-page
document that matched fell outside the three-page rule. Both prompts presuppose
a market-definition section, so for a decision without market analysis the
model produced precedent-style passages, and eight records reproduce the
simplified-procedure recital itself. Version 1.1.0 matches the phrases
independently of whitespace and applies the rule to documents of at most three
pages, which excludes all nine short-form documents of the pilot. The fidelity
script removes absent records and, optionally, all records from documents of at
most three pages, which produced the cleaned release.

\paragraph{Prompt examples.} The four example passages in the extraction prompt
appear in 30 records across 11 cases that are not their source. The structural
test accepts them, and the fidelity test flags all 30. The next version of the
prompt omits the examples.

\paragraph{Boundaries, duplicates and wording.} Near-verbatim records show the
model splicing across footnotes and occasionally omitting a paragraph. At the
other end, 261 records are shorter than 26 words, mostly single conclusion
sentences, and 14 exceed 668 words. In addition, 117 records repeat the text of
another record, mostly one passage stored under two topic labels. These
segmentation questions are part of the gold-standard study. Most remaining
structural failures are genuine definitions in wording the 18 markers do not
cover, notably ``national in scope''.

% -----------------------------------------------------------------------------
\section{Limitations and Future Work}
\label{sec:limitations}
% -----------------------------------------------------------------------------

\emph{Gold standard.} Precision, recall and F1 against an independently
annotated set of all definitions in a decision remain to be measured. The
fidelity test establishes that 88\% of records are copied from their source;
whether they are the right passages and whether any were missed is the subject
of a planned annotation of ten decisions, for which
\texttt{make\_annotation\_sheet.py} provides the sheet and protocol.
\emph{Baselines.} The two-stage design and the choice of Gemini 2.0 Flash,
taken for its free-tier rate limits and long context, have yet to be compared
with a single-call variant, a heading-based rule baseline and other LLMs. The
first two are implemented in \texttt{baselines/} and feed the same evaluation
scripts. \emph{Annotation.} The manual inspection and the reference set were
each labelled once, without a second annotator. \emph{Scope.} The pipeline
processes English-language decisions of one authority, and its output depends
on the quality of the PDF text.

% -----------------------------------------------------------------------------
\section{Conclusion}
\label{sec:conclusion}
% -----------------------------------------------------------------------------

Lextract turns heterogeneous Commission decisions into 4,105 structured
market-definition records with a two-call LLM pipeline whose prompts, settings
and filters are documented in full. Across the collection, 79\% of records are
verbatim copies of their source and a further 9\% are near-verbatim. The
source-fidelity test identifies the remaining records, most of them from short
decisions without market analysis, which a resemblance-based check accepts, and
produces a cleaned release of 3,607 records. Combining LLM extraction with a
deterministic check against the source text is inexpensive and applies to any
extraction pipeline over legal documents.

\paragraph{Acknowledgements.}
The author is grateful to Professor Thibault Schrepel of Stanford Law School for
his advisement and guidance, and to the JURIX 2026 reviewers. This research
received no funding from any government agency, university, company, or
non-profit organization.

% -----------------------------------------------------------------------------
% References
% -----------------------------------------------------------------------------
\bibliographystyle{vancouver}
\bibliography{paper}

\end{document}
